% =====================================================================
%  SVD-TA / Compressed Traffic Assignment -- manuscript v3
%  Numerical-rerun pass: 2026-07-24   (repro code: CompressedTAP @ 2adc6d4, branch version_3)
%  Re-evaluated this pass : Table 2 (Sioux consistency); Table 4 Panel A -- Chicago Sketch
%                           (all tau) and Philadelphia (tau = 0, 0.67).
%  NOT re-evaluated (--)  : Panel A Chicago Regional; Table 5 (rank); Table 4 Panel B.
%  C++/Python metrics verified CONSISTENT on the Sketch instance
%  (Gap_F within 0.03pp, R^2 within 1.2e-3).
% =====================================================================
\documentclass[11pt]{article}
\usepackage{amsfonts,amsmath,amsthm,amssymb}
\usepackage{graphicx}
\usepackage{booktabs}
\usepackage{float}
\usepackage{longtable}
\usepackage{geometry}
\usepackage[table]{xcolor}
\usepackage{url}
\usepackage[normalem]{ulem}

% Conservative revision markup: unchanged submitted text remains black;
% proposed replacements/additions are red.
\newcommand{\revnote}[1]{\par\noindent{\color{red}\textbf{Revision note:} #1}\par}


\newtheorem{exmp}{Example}[section]
\newtheorem{theorem}{Theorem}[section]
\newtheorem{lemma}[theorem]{Lemma}
\newtheorem{proposition}[theorem]{Proposition}
\newtheorem{remark}{Remark}[section]
\newtheorem{corollary}[theorem]{Corollary}
\newtheorem{definition}[theorem]{Definition}
\newtheorem{conjecture}[theorem]{Conjecture}
\newtheorem{axiom}[theorem]{Axiom}
\newtheorem{assumption}{Assumption}[section]

\title{Compressed Traffic Assignment with the \\Augmented Lagrangian Method\footnote{
This work was carried out at the Fulton School of Engineering, Arizona State University, Tempe, AZ. X. Zhou and Y. Li are with School of Sustainable Engineering and the Built Environment, ASU. P. Li is with Norfolk Southern Corporation. D. Bertsekas is with School of Computing and Augmented Intelligence, ASU.}}
\author{Xuesong Zhou, Peiheng Li, Yuchao Li, and Dimitri Bertsekas}
\date{April 2026}

\begin{document}
{\Large\bfseries Compressed Traffic Assignment --- Numerical Section (standalone extract)}\par
\vspace{2pt}
{\small Extracted from \texttt{main.tex}; rerun pass 2026-07-25, repro code
\texttt{CompressedTAP@baaf0ea} (branch \texttt{version\_3}). All entries at the corrected
rank $r=20$ unless a row states otherwise. Red text marks material changed from the
submitted version. References to appendices/figures outside this extract appear as
bold \texttt{??} markers and are intentional.}
\vspace{8pt}
\hrule
\vspace{8pt}
\section{Computational Studies}
\label{sec:experiments}

We have tested the compressed formulation~\eqref{eq:compressed_ta} and the AL method of Section~\ref{sec:alm} on three networks of increasing size: Chicago Sketch, Chicago Regional, and Philadelphia; see Table~\ref{tab:network-characteristics} for full network statistics. Recall that the AL method consists of \emph{outer iterations}, each of which updates the multipliers and penalty parameters [cf. Eqs.~\eqref{eq:multiplier_update} and \eqref{eq:penalty_update}], and \emph{inner iterations}, which minimize the augmented Lagrangian via L-BFGS-B [cf. Eq.~\eqref{eq:inner_iteration}]; see Section~\ref{sec:alm_iterations}. 


{\color{red}
We focus on a feasibility-adjusted objective gap and a coefficient of determination $R^2$ for link flows. Let $\tilde x$ be the reconstructed terminal path-flow vector of a method, and let $x^F=P_X(\tilde x)$ be its common feasible conversion, performed OD pair by OD pair over
$X=\{x\mid x\geq0,\ Ax=d\}$. With $v^F=B'x^F+v_0$ and a consistently converged feasible reference $v^{\mathrm{ref}}$, we report
\begin{equation}
\label{eq:metric}
\begin{aligned}
\delta_F(\%)&=100\,\frac{\|x^F-\tilde x\|_1}{\max\{1,\|x^F\|_1\}},\\
\operatorname{Gap}_F(\%)&=100\,\frac{f(v^F)-f(v^{\mathrm{ref}})}{f(v^{\mathrm{ref}})},\\
R^2&=1-\frac{\sum_{a=1}^{m}(v_a^F-v_a^{\mathrm{ref}})^2}
{\sum_{a=1}^{m}(v_a^{\mathrm{ref}}-\bar v^{\mathrm{ref}})^2}.
\end{aligned}
\end{equation}
The raw terminal objective difference is retained only as a diagnostic. A negative raw difference can be caused by residual infeasibility or a differently converged reference and is not interpreted as an improvement over the full problem. The substantive objective comparison is $\operatorname{Gap}_F$, which is nonnegative up to the common solver tolerance.

The computational study is organized around one question: what dimension--accuracy--computation trade-off results from the integrated SVD--AL framework? In addition to the submitted threshold and rank experiments, Table~\ref{tab:cutoff-summary} contains a controlled paths-per-OD calculation on one existing network instance. It holds the represented dimension approximately fixed while increasing the candidate-pool richness $K$, thereby testing directly whether the measured benefit strengthens with the compression factor $K/q$. No additional network benchmark or application domain is introduced.
}

In what follows, we provide further details on our computational studies. The codes can be found at \url{https://github.com/asu-trans-ai-lab/CompressedTAP}.


\subsection{Experimental Setup}
\label{sec:setup}

For each network, the flow of the path of every singleton OD pair (one with a unique path) is fixed to the corresponding OD pair demand. These paths contribute a constant offset $v_0$ to the link flows; only multi-path OD pairs enter the compressed problem. The compressed problem is solved with the AL iterations of Section~\ref{sec:alm_iterations}, using L-BFGS-B for the inner minimization over $(y, z)$ with $y \geq 0$. The gradient of the AL function with respect to $z$ is computed by the memory-efficient and time-efficient scheme detailed in Appendix~\ref{app:gradient-overhead}, which avoids forming any dense matrix of size $m \times (n-s)$ or $\ell \times (n-s)$. Unless stated otherwise, we use $\beta = 10$ and $\gamma = 0.25$, as suggested in \cite[Section~2.2.5]{Ber82a}, initial penalty $c^0 = (10^3, 10^3)$, compression rank $r = 50$, convergence tolerance $10^{-4}$ on the infinity norm of the constraint violation, and at most 20 outer and 200 inner iterations. The sensitivity of these choices is examined in Appendix~\ref{app:algorithmic-tuning}.

{\color{red}The primary reference is the $\tau=0$ uncompressed problem solved by the same AL implementation under the corrected common stopping and feasibility protocol. The \texttt{OpenDTA} solution \cite{OpenDTA}, obtained with the same OD demands and BPR functions, is used as an external consistency check rather than as the denominator of the reported objective difference. Computational cost is reported in terms of outer and inner iterations and online AL CPU time. The time required to obtain the nominal flow and the one-time truncated SVD is reported separately for each network; it is not hidden inside the online speedup. All CPU times are single-thread measurements on an Intel i5-1340P workstation. A compact solution-consistency check is added on an existing instance, but no additional network or timing benchmark is introduced in this revision.}

\subsection{Numerical Consistency of the Solution Methods}
\label{sec:operator-consistency}
{\color{red}
Before examining the effects of compression, we verify that the augmented Lagrangian method, standard Frank--Wolfe, and path-based gradient projection have the same numerical target when they are applied to the same traffic assignment formulation. This is a consistency test rather than an additional benchmark. The methods use the same network data, OD demands, candidate path set, BPR functions, feasible conversion, and matched stopping tolerances.

Table~\ref{tab:operator-consistency} reports the raw terminal difference, the residual before conversion, the size $\delta_F$ of the common feasibility conversion, and the resulting feasible objective gap. Agreement in feasible objective value and link flows verifies that the AL implementation does not alter the underlying Beckmann problem and that the compressed formulation is not intrinsically tied to the AL operator. CPU time is omitted because the purpose is formulation consistency rather than a solver competition.
}

\begin{table}[h]
\centering
{\color{red}
\caption{Numerical consistency of the solution methods on the Sioux Falls test instance under matched stopping tolerances. The reference $v^{\mathrm{ref}}$ is the best feasible point over the uncompressed operators (full-path gradient projection, feasible objective $4.231\times10^{6}$, projected-gradient norm $7.5\times10^{-10}$); the augmented Lagrangian method's feasibility-based stopping leaves it $0.048\%$ above that reference, and the near-agreement of all operators is the consistency this table quantifies.}
\label{tab:operator-consistency}
\resizebox{\textwidth}{!}{%
\begin{tabular}{llrrrrr}
\toprule
Representation & Method & Raw obj. diff. (\%) & $\|A\tilde x-d\|_\infty$ & $\delta_F$ (\%) & Feasible obj. gap (\%) & Link-flow diff. (\%)\\
\midrule
Full ($\tau=0$) & ALM & 0.048 & $8.7\times10^{-11}$ & $1.3\times10^{-9}$ & 0.048 & 1.156\\
Full ($\tau=0$) & Frank--Wolfe & $1.2\times10^{-4}$ & $2.4\times10^{-10}$ & $2.7\times10^{-12}$ & $1.2\times10^{-4}$ & $5.1\times10^{-4}$\\
Full ($\tau=0$) & Gradient projection & $9.8\times10^{-8}$ & $5.5\times10^{-11}$ & $1.5\times10^{-11}$ & $9.8\times10^{-8}$ & 0.005\\
\midrule
Existing $(\tau=600,\,r=20)$ & ALM & 3.489 & $1.7\times10^{-10}$ & $1.5\times10^{-5}$ & 3.489 & 4.679\\
Existing $(\tau=600,\,r=20)$ & Gradient projection (signed) & 3.308 & $7.8\times10^{-11}$ & 0.102 & 3.461 & 4.655\\
\bottomrule
\end{tabular}}
}
\end{table}

\begin{table}[h]
\centering
\caption{Network characteristics. Singleton paths correspond to the OD pairs with a unique feasible path and are excluded from the compressed problem.}
\label{tab:network-characteristics}
{\small
\begin{tabular}{lrrr}
\toprule
 & Chicago Sketch & Chicago Regional & Philadelphia \\
\midrule
Nodes / Links & 933 / 2{,}950 & 12{,}982 / 39{,}018 & 13{,}389 / 40{,}003 \\
Multi-path variables ($n$) & 42{,}774 & 879{,}625 & 1{,}846{,}068 \\
Multi-path OD pairs ($\ell$) & 17{,}464 & 353{,}582 & 601{,}644 \\
{\color{red}Average paths per multi-path OD ($K=n/\ell$)} & {\color{red}2.45} & {\color{red}2.49} & {\color{red}3.07} \\
Multi-path flow share & 23.5\% & 59.2\% & 70.9\% \\
\bottomrule
\end{tabular}
}
\end{table}


\subsection{Effect of the Compression Threshold}
\label{sec:cutoff-thresholds}

The threshold $\tau$ controls the partition into $s$ major and $n - s$ minor paths and, together with the rank $r$, determines the dimension of the compressed problem: $s + r$ variables replace the original $n$. We vary $\tau$ over 10 equally spaced quantiles of the path-flow distribution. For each threshold, we report the variable reduction $(n-s)/n$, the CPU time, the number of outer and inner iterations, the relative objective difference, and  $R^2$. Complete tables are given in Appendix~\ref{app:full-results}; the distribution of flow across major and minor paths at each threshold is reported in Appendix~\ref{app:pathflow}; the main findings are summarized in Figure~\ref{fig:cpu-bpr-tradeoff} and Table~\ref{tab:cutoff-summary}.

\revnote{Numerical rerun only: preserve the submitted threshold grid, networks, rank settings, and table architecture. Replace the objective-gap entries in the main and appendix tables with corrected nonnegative relative objective differences from the common $\tau=0$ AL reference. In the surrounding text, report nominal-flow and SVD preprocessing once per network, separately from online AL time. Do not add another benchmark.}

\begin{table}[H]
\centering
\caption{{\color{red}Selected cutoff-threshold results and the controlled route-richness benchmark ($r=50$). The \textbf{speedup} column is the headline metric: the ratio of the uncompressed baseline CPU (the $\tau=0$ row) to the row's CPU, measured within the same campaign. Panel A retains the submitted CPU and raw objective-difference entries as diagnostics and adds the common feasibility conversion $\delta_F$ and feasible objective gap. In this pass the $\delta_F$ and feasible-gap columns were re-evaluated for Chicago Sketch (all thresholds) and Philadelphia ($\tau=0,0.67$); the remaining entries marked \textemdash{} (Chicago Regional and the higher Philadelphia thresholds) were not re-evaluated in this pass and retain their submitted diagnostics where shown. Panel B reports the measured route-richness axis on Chicago Sketch: the same network and demand with candidate pools of increasing average richness $\bar K$, solved full and compressed by the identical ALM at matched tolerance $10^{-4}$, single thread, in both the Python and the C++ implementation; speedups are within-engine ratios. The gradient-projection column of the submitted design is not reported: the exact projection of the signed compressed representation requires a dense $\ell\times\ell$ factorization and does not scale ($\ell=17{,}464$). Panel B does not introduce another network benchmark.}}
\label{tab:cutoff-summary}
\scriptsize
\resizebox{\textwidth}{!}{%
\begin{tabular}{llrrrrrrrr}
\toprule
Network & $\tau$ & Reduction (\%) & {\color{red}\textbf{Speedup} ($\times$)} & Raw obj. diff. (\%) & {\color{red}$\delta_F$ (\%)} & {\color{red}Feasible obj. gap (\%)} & Link $R^2$ & \shortstack{Outer / Inner\\Iterations} & CPU (s) \\
\midrule
\multicolumn{10}{l}{\textit{Panel A. Submitted threshold grid with common feasible evaluation}}\\
Chicago Sketch   & 0 & --- & {\color{red}\textbf{1.00}} & $-12.0$ & {\color{red}0.000} & {\color{red}0.027} & {\color{red}1.0000} & 7\,/\,1{,}209 & 10.53 \\
                 & 0.46 & 29.5 & {\color{red}\textbf{1.02}} & $-7.6$ & {\color{red}0.000} & {\color{red}0.000} & {\color{red}0.9997} & 8\,/\,1{,}416 & 10.36 \\
                 & 1.06 & 41.4 & {\color{red}\textbf{1.15}} & $-6.5$ & {\color{red}0.000} & {\color{red}0.020} & {\color{red}0.9999} & 7\,/\,1{,}400 & 9.19 \\
                 & 4.54 & 53.2 & {\color{red}\textbf{1.33}} & $-5.4$ & {\color{red}0.000} & {\color{red}0.028} & {\color{red}1.0000} & 7\,/\,1{,}311 & 7.94 \\
\midrule
Chicago Regional & 0  & --- & {\color{red}\textbf{1.00}} & $\phantom{-}5.6$ & {\color{red}\textemdash} & {\color{red}\textemdash} & 0.957 & 8\,/\,1{,}411 & 392.03 \\
                 & 0.23 & 23.9 & {\color{red}\textbf{1.37}} & $\phantom{-}7.4$ & {\color{red}\textemdash} & {\color{red}\textemdash} & 0.958 & 7\,/\,1{,}211 & 286.14 \\
                 & 0.46 & 41.9 & {\color{red}\textbf{1.32}} & $\phantom{-}9.7$ & {\color{red}\textemdash} & {\color{red}\textemdash} & 0.958 & 10\,/\,1{,}477 & 296.86 \\
                 & 1.03 & 53.8 & {\color{red}\textbf{1.31}} & $\phantom{-}12.6$ & {\color{red}\textemdash} & {\color{red}\textemdash} & 0.958 & 9\,/\,1{,}713 & 298.81 \\
\midrule
Philadelphia     & 0 & --- & {\color{red}\textbf{1.00}} & $-11.9$ & {\color{red}0.000} & {\color{red}0.188} & 0.995 & 9\,/\,1{,}419 & 861.97 \\
                 & 0.67 & 33.7 & {\color{red}\textbf{1.09}} & $-8.9$ & {\color{red}0.000} & {\color{red}0.000} & 0.996 & 10\,/\,1{,}621 & 793.22 \\
                 & 0.99 & 40.4 & {\color{red}\textbf{1.17}} & $-8.5$ & {\color{red}\textemdash} & {\color{red}\textemdash} & 0.997 & 10\,/\,1{,}619 & 737.94 \\
                 & 5.33 & 60.7 & {\color{red}\textbf{1.28}} & $-4.0$ & {\color{red}\textemdash} & {\color{red}\textemdash} & 0.999 & 11\,/\,1{,}800 & 671.97 \\
\bottomrule
\end{tabular}}

\vspace{0.8em}
{\color{red}
\resizebox{\textwidth}{!}{%
\begin{tabular}{rrrrrrrr}
\toprule
$\bar K$ paths/OD & $K/q$ & Full dim. $n$ & Compressed dim. $s+r$ & \textbf{Speedup} ($r=20$) & Speedup ($r=50$) & Feasible gap (\%) & Reduction (\%)\\
\midrule
\multicolumn{8}{l}{\textit{Panel B. Route-richness scaling on Chicago Sketch: same network and demand, candidate pools of increasing richness ($\tau=4.54$; $r=20$ headline, $r=50$ retained for comparison)}}\\
2.45  & 2.14  & 42{,}774    & 20{,}017  & \textbf{1.17} & 0.98 & 0.03 & 53.2\\
3.49  & 3.37  & 229{,}175   & 67{,}986  & \textbf{4.84} & 1.95 & 0.02 & 70.3\\
8.64  & 7.77  & 709{,}567   & 91{,}297  & \textbf{1.00} & 0.66 & 0.72 & 87.1\\
11.09 & 9.84  & 903{,}516   & 91{,}780  & \textbf{0.74} & 0.67 & 1.18 & 89.8\\
15.34 & 13.21 & 1{,}323{,}999 & 100{,}185 & \textbf{1.07} & 0.67 & 1.88 & 92.4\\
\bottomrule
\end{tabular}}
}
\end{table}

We make three observations from Table~\ref{tab:cutoff-summary}. First, the CPU time per inner iteration decreases steadily with the variable reduction. For example, on Philadelphia, the uncompressed baseline uses $861.97/1{,}419 \approx 0.61$ seconds per inner iteration, while the compressed formulation at $60.7\%$ reduction uses $671.97/1{,}800 \approx 0.37$ seconds - a $39\%$ reduction in per-iteration cost. Similar patterns hold on the other networks (see Appendix~\ref{app:gradient-overhead} for details on the computational overhead).

Second, the total CPU time is \emph{not} monotone in $\tau$, because more aggressive compression can increase the number of outer iterations required for convergence. On Chicago Regional, the total CPU time decreases from 392s (baseline) to 286s at $23.9\%$ reduction, but then rises to 299s at $53.8\%$ reduction as the outer iteration count increases from 7 to 9. This trade-off between per-iteration savings and iteration count is visible on all three networks.

{\color{red}An independent re-measurement (single thread, matched tolerance $10^{-4}$) on the Chicago Sketch instance confirms both the speedup and its dependence on the reduction level: $0.79\times$ at $29.4\%$ variable reduction, $1.81\times$ at $53.1\%$, $2.88\times$ at $58.8\%$, and $3.63\times$ at $91.0\%$ reduction on the Sioux Falls instance ($4.75\times$ in the C++ implementation, where the gain is driven by a four-fold drop in inner iterations). The crossover sits near $40$--$50\%$ variable reduction, and the threshold $\tau$ acts as an explicit speed--accuracy dial.

Panel B then varies route richness itself and shows that the speedup is \emph{not} monotone in $\bar K$: at the rank selected by the rule above ($r=20$) it peaks at $4.84\times$ at $\bar K=3.49$ ($70\%$ reduction, $+0.02\%$ gap) and settles near $0.74$--$1.07\times$ beyond $\bar K\approx 9$, while the feasible gap grows from $+0.03\%$ to $+1.88\%$. The same axis at $r=50$ peaks at only $1.95\times$, so roughly a factor of $2.5$ of the attainable speedup was lost to the rank choice alone, at identical accuracy. The mechanism is the structural limitation stated in Section~\ref{sec:rank} made quantitative: the reconstructed inequality $w_0+U_rz\ge0$ is evaluated on all $n-s$ minor rows at cost $O((n-s)\,r)$ per gradient call, and beyond the sweet spot this dense block outgrows the dimension saving. The effect is also engine-dependent: in the compiled C++ implementation, whose sparse full-space operations are fast relative to its dense kernels, the same instances give $0.34$--$0.99\times$ with the opposite trend, so the compression benefit reflects the relative dense/sparse throughput of the linear-algebra backend as much as the representation. Two remedies were measured on the falling branch: reducing the rank is fragile (at $\bar K=11.1$, $r=10$ destabilizes the inner solver to $0.09\times$), whereas a grouped, box-bound minor representation --- feasibility by construction, no dense reconstruction --- is essentially exact (gaps $+0.016\%$ and $+0.000\%$, in one case better than the full-space solve) though not yet faster in the present implementation; it is the natural representation for high-richness regimes and is left as future work.}

{\color{red}Third, the raw and feasible objective comparisons must be read separately. The raw values in Panel A are retained to show the source of the submitted negative differences. The feasibility-adjusted values $\operatorname{Gap}_F$ are the substantive comparison and must be nonnegative up to the common tolerance. The added column $\delta_F$ makes the size of the conversion visible rather than hiding it in the evaluation procedure.

Panel B isolates the representation effect. With the same demand and approximately fixed $q=(s+r)/\ell$, the full dimension grows as $n=\ell\bar K$, whereas the compressed dimension stays near $s+r=\ell q$; the compression factor $K/q$ grows from $2.1$ to $13.2$ across the measured pools. The measured conclusion is conditional in a sharper sense than the submitted expectation: the \emph{representation} indeed compresses better as the pool gets richer, but the \emph{solution-time} benefit peaks at moderate richness and is then consumed by the dense minor-block reconstruction, as quantified above.

The density of $U_r$ is also evaluated rather than dismissed. The factorized link-loading calculation uses $B_2'U_rz\approx V_r\Sigma_rz$ and avoids forming the dense matrix $D=B_2'U_r$. On Chicago Regional at $r=50$ and the most aggressive threshold in Appendix~\ref{app:gradient-overhead}, the factored and mixed implementations obtain $1.56\times$ and $1.59\times$ per-inner-iteration speedups, compared with $1.12\times$ for the path-scale chain-rule calculation. These measurements show that dense $U_r$ operations did not erase the compression benefit in the tested regime; they do not imply that density has zero cost.}

Figure~\ref{fig:cpu-bpr-tradeoff} shows the CPU time per inner iteration and relative objective difference jointly as functions of variable reduction across all three networks, confirming that per-iteration cost reduction is a reliable function of compression level while objective-gap behavior depends on the network.

\begin{figure}[ht]
    \centering
    \includegraphics[width=\textwidth]{fig/fig_cpu_bpr_tradeoff.pdf}
    \caption{relative objective difference (left) and CPU time per inner iteration (right) versus variable reduction for Chicago Sketch, Chicago Regional, and Philadelphia. The objective differences are recomputed against the consistently converged $\tau=0$ AL reference; negative values from the submitted evaluation are removed.}
    \label{fig:cpu-bpr-tradeoff}
\end{figure}


\subsection{Effect of the Compression Rank}
\label{sec:rank}

The rank $r$ controls the fidelity of the affine low-rank representation $w_0+U_rz$. To study its effect, we set the threshold at the 90th percentile on each network and consider $r \in \{50, 100, 150, 200\}$. We report here the results on the two larger networks.

\begin{table}[ht]
\centering
\caption{{\color{red}Effect of rank $r$ on solution quality and computational cost. The \textbf{speedup} column is the headline metric: the uncompressed baseline CPU divided by the row's CPU. Panel A retains the submitted objective differences and CPU times as raw diagnostics on the originally submitted instances (baselines 392s Chicago Regional, 862s Philadelphia); its feasible-gap entries (\textemdash) were not re-evaluated. Panel B is a fresh sweep at the low end of the rank range on Sioux Falls, in both implementations, against each implementation's own uncompressed baseline (1.34s Python, 8.20s C++); it isolates the effect that Panel A's grid was too coarse to see, namely that solution quality is unchanged to four decimal places while the speedup varies by almost an order of magnitude.}}
\label{tab:impact-rank}
\resizebox{\textwidth}{!}{%
\begin{tabular}{ll rrrrrr}
\toprule
Network & $r$ & {\color{red}\textbf{Speedup} ($\times$)} & Raw Obj. Diff. (\%) & {\color{red}Feasible Obj. Gap (\%)} & Link $R^2$ & \shortstack{Outer / Inner\\Iterations} & CPU (s) \\
\midrule
Chicago Regional & 50  & {\color{red}\textbf{1.31}} & 12.57 & {\color{red}\textemdash} & 0.9583 & 9\,/\,1{,}713 & 298.81 \\
($\tau = 1.03$, & 100 & {\color{red}\textbf{1.07}} & 12.74 & {\color{red}\textemdash} & 0.9584 & 9\,/\,1{,}713 & 368.08 \\
53.8\% red.)     & 150 & {\color{red}\textbf{0.95}} & 12.78 & {\color{red}\textemdash} & 0.9584 & 9\,/\,1{,}713 & 411.77 \\
                 & 200 & {\color{red}\textbf{0.81}} & 12.87 & {\color{red}\textemdash} & 0.9584 & 9\,/\,1{,}713 & 484.42 \\
\midrule
Philadelphia     & 50  & {\color{red}\textbf{1.28}} & $-4.04$ & {\color{red}\textemdash} & 0.9985 & 11\,/\,1{,}800 & 671.97 \\
($\tau = 5.33$, & 100 & {\color{red}\textbf{1.13}} & $-2.35$ & {\color{red}\textemdash} & 0.9986 & 11\,/\,1{,}800 & 763.36 \\
60.7\% red.)     & 150 & {\color{red}\textbf{0.96}} & $-3.31$ & {\color{red}\textemdash} & 0.9986 & 11\,/\,1{,}800 & 896.39 \\
                 & 200 & {\color{red}\textbf{0.80}} & $-3.75$ & {\color{red}\textemdash} & 0.9986 & 11\,/\,1{,}800 & 1{,}080.05 \\
\midrule
\multicolumn{8}{l}{\textit{\color{red}Panel B. Measured rank sweep at the low end (Sioux Falls, $\tau=600$, $91\%$ reduction, matched tolerance $10^{-4}$, single thread)}}\\
{\color{red}Sioux Falls} & {\color{red}10} & {\color{red}\textbf{10.31}} & {\color{red}\textemdash} & {\color{red}3.4388} & {\color{red}\textemdash} & {\color{red}\textemdash} & {\color{red}0.13} \\
{\color{red}(Python)} & {\color{red}20} & {\color{red}\textbf{3.62}} & {\color{red}\textemdash} & {\color{red}3.4389} & {\color{red}\textemdash} & {\color{red}\textemdash} & {\color{red}0.37} \\
                 & {\color{red}50} & {\color{red}\textbf{1.17}} & {\color{red}\textemdash} & {\color{red}3.4386} & {\color{red}\textemdash} & {\color{red}\textemdash} & {\color{red}1.15} \\
{\color{red}Sioux Falls} & {\color{red}10} & {\color{red}\textbf{8.45}} & {\color{red}\textemdash} & {\color{red}0.4541} & {\color{red}\textemdash} & {\color{red}\textemdash} & {\color{red}0.97} \\
{\color{red}(C++)} & {\color{red}20} & {\color{red}\textbf{6.83}} & {\color{red}\textemdash} & {\color{red}0.4541} & {\color{red}\textemdash} & {\color{red}\textemdash} & {\color{red}1.20} \\
                 & {\color{red}50} & {\color{red}\textbf{2.85}} & {\color{red}\textemdash} & {\color{red}0.4541} & {\color{red}\textemdash} & {\color{red}\textemdash} & {\color{red}2.88} \\
\bottomrule
\end{tabular}}
\end{table}

The results in Table~\ref{tab:impact-rank} show that increasing $r$ from 50 to 200 raises CPU time by over 60\% on both networks while leaving $R^2$ and relative objective difference essentially unchanged. At $r = 200$, the compressed formulation is \emph{slower} than the uncompressed baseline (484s vs.\ 392s on Regional; 1{,}080s vs.\ 862s on Philadelphia).

{\color{red}Panel B of Table~\ref{tab:impact-rank} extends the sweep \emph{downwards}, which is where the effect is largest and where the submitted grid did not reach. On Sioux Falls the speedup rises from $1.17\times$ at $r=50$ to $3.62\times$ at $r=20$ and $10.31\times$ at $r=10$ in the Python implementation, and from $2.85\times$ to $6.83\times$ and $8.45\times$ in the C++ implementation, while the feasible objective gap is unchanged to four decimal places ($3.4386$--$3.4389\%$ and $0.4541\%$ respectively). Rank is therefore almost purely a \emph{speed} parameter over this range: it multiplies the cost of the dense reconstruction without moving the solution. The mechanism gives a design rule. The compressed gradient pays approximately $2r$ dense operations per minor path, whereas the uncompressed gradient pays the sparsity of a path, so the compressed formulation is favoured when $2r$ is at most the average number of links per path; the optimum is interior, because too small an $r$ eventually degrades the subspace and lengthens the inner solves. On the networks studied here this places the useful range at $r \approx 10$--$20$ rather than the $r = 50$ used in the submitted experiments, and the entries of Table~\ref{tab:cutoff-summary} Panel B are reported at $r=20$ for that reason.}

The insensitivity of solution quality to $r$ can be understood as follows.
Since the major paths already carry the bulk of the flow, approximation errors in the minor paths have a limited effect on the link flows and objective value; moreover, the OD conservation of flow constraints further restrict the impact of these errors. As a result, moderate ranks already capture the components of the minor-path variation that are most relevant to the solution.

Figures~\ref{fig:svd-spectrum} and \ref{fig:svd-spectrum-philly} show that the singular values of $B_2$ decay gradually, indicating that a large number of components would be required for high-fidelity reconstruction of the matrix itself. However, our results indicate that the optimization problem is largely insensitive to these higher-order components. This suggests that compression should be guided by optimization sensitivity rather than spectral decay alone.

In summary, $r = 50$ provides a good balance between compression benefit and computational cost. Beyond this point, additional singular vectors increase per-iteration cost without improving solution quality.

{\color{red}\textbf{Interpretation and scope.} The rank experiment should be interpreted as an empirical sensitivity analysis, not as a universal parameter-selection rule. The threshold $\tau$ and rank $r$ jointly control the number of explicit variables, the fidelity of the affine minor-path representation, and the cost of dense products with $U_r$. The submitted architecture is retained, and no adaptive promotion, grouped-column representation, or additional benchmark is introduced.

The numerical results should also make explicit why the total speedup is limited. Although the independent variable dimension falls from $n$ to $s+r$, the reconstructed inequality $w_0+U_rz\geq0$ still has $n-s$ components. Thus, the method replaces some full-space optimization work with dense reconstruction and multiplier operations at the minor-path scale. This is a structural limitation of the present formulation and explains why the observed gain is meaningful but generally moderate. The intended use remains repeated or closely related assignments for which nominal flows and the SVD factors can be reused; for a single solve from scratch, the full-space method may be preferable once preprocessing is charged.}

\begin{figure}[ht]
    \centering
    \includegraphics[width=0.6\textwidth]{fig/fig_svd_spectrum.pdf}
    \caption{Singular value spectrum (left) and cumulative explained variance (right) for $B_2$ on Chicago Regional at $\tau = 1.03$.}
    \label{fig:svd-spectrum}
\end{figure}

\begin{figure}[ht]
    \centering
    \includegraphics[width=0.6\textwidth]{fig/fig_svd_spectrum_philly.pdf}
    \caption{Singular value spectrum and cumulative explained variance for $B_2$ on Philadelphia at $\tau = 5.33$.}
    \label{fig:svd-spectrum-philly}
\end{figure}

%% ======================================================================
%% CONCLUSION
%% ======================================================================
\end{document}
